% Generated draft: requires mathematical and source review.
\documentclass[11pt]{article}
\usepackage[margin=25mm]{geometry}
\usepackage{fontspec}
\setmainfont{Latin Modern Roman}
\usepackage{amsmath,amssymb,mathtools,unicode-math}
\setmathfont{Latin Modern Math}
\usepackage{enumitem,needspace,xurl,hyperref}
\setlength{\parindent}{0pt}
\setlength{\parskip}{5pt}
\setlength{\emergencystretch}{3em}
\newcommand{\N}{\mathbb{N}}
\newcommand{\Z}{\mathbb{Z}}
\newcommand{\Q}{\mathbb{Q}}
\newcommand{\R}{\mathbb{R}}
\newcommand{\C}{\mathbb{C}}
\newcommand{\F}{\mathbb{F}}
\newcommand{\A}{\mathbb{A}}
\newcommand{\PP}{\mathbb{P}}
\newcommand{\E}{\mathbb{E}}
\newcommand{\eps}{\varepsilon}
\newcommand{\dd}{\,\mathrm{d}}
\begin{document}
\section*{UnsolvedMath ID: 8800079}
\textbf{Source problem code:} catalogue entry
\subsection{Definitions and mathematical statement}
Construct an efficient public-key encryption scheme analogous to Regev's LWE cryptosystem, with security based on specified worst-case problems on ideal lattices through an appropriate quantum reduction. ``Quantum-SVP-based'' concerns the reduction/hardness basis and does not itself require quantum encryption or decryption algorithms.

A scheme consists of $\operatorname{KeyGen}(1^\lambda)$, $\operatorname{Enc}_{pk}(m)$, and $\operatorname{Dec}_{sk}(c)$, with correctness except for any explicitly bounded negligible failure probability. A standard security goal is indistinguishability under chosen-plaintext attack, with efficient adversarial advantage negligible in $\lambda$. The reduction and adversary models must be stated.

For context, a ring-LWE-type sample in a chosen ring quotient $R_q=R/qR$ has the schematic form
\[
(a,b)=(a,as+e)\in R_q^2,
\]
where $a$ is uniform and the secret and error distributions are specified. Canonical versus coefficient embeddings, dual-ring variants, field/ring family, modulus, and noise parameters matter for the corresponding ideal-lattice theorem. The source does not fix them or an approximation factor, so no universal reduction for arbitrary rings/noise is inferred.

The supplied summary attributes a solution to ring-LWE constructions, especially Lyubashevsky--Peikert--Regev (2010), with earlier related 2009 work. This is a reported existence resolution, not an unconditional proof of hardness. Its Kyber mention is an application connection: a module-lattice KEM is not silently identified with a rank-one ideal-lattice encryption theorem.
\subsection{Short English statement}
Build an efficient ideal/ring-lattice analogue of Regev encryption, with an explicit quantum security reduction and parameter assumptions. The supplied summary reports ring-LWE constructions as resolving the existence problem.
\subsection{Sources}
\begin{itemize}
\item Daniele Micciancio, IPAM \emph{Workshop on Number Theory and Cryptography --- Open Problems} (2006), attachment 6240, PDF page 41, schedule source \url{https://www.ipam.ucla.edu/programs/workshops/workshop-i-number-theory-and-cryptography-open-problems/?tab=schedule}; accessed 2026-07-29. The supplied provenance records a PDF/legacy-PPT visual audit.
\item Stehl\'e--Steinfeld--Tanaka--Xagawa (2009), and Lyubashevsky--Peikert--Regev (2010), ``On ideal lattices and learning with errors over rings,'' as supplied resolution attributions. The summary also mentions Kyber without a separate publication citation.
\item Local registry \texttt{data/unsolvedmath/problems.json}, permanent UnsolvedMath ID $8800079$, source code catalogue entry, \emph{Ideal-lattice Regev cryptosystem}. Dataset attribution: Ulam AI and the UnsolvedMath Contributors.
\end{itemize}
\end{document}
